% Caption for output/publication/final/index_planes_sfh_families.{pdf,png}
\caption{%
The AGB-on index planes of Fig.~\ref{fig:index_planes_agb_on_off} for three
alternative forms of the post-quenching star formation rate. All families share the
same delayed-$\tau$ rise and the same 2000 draws of $(t_q, \tau_q, \log Z)$; they differ
only after $t_q$. Top row: a linear ramp to zero over $2 \ln 2\, \tau_q$, the FSPS
${\tt sfh} = 5$ form with the same star-formation half-life as the exponential decline.
Middle row: an abrupt truncation, the FSPS ${\tt sfh} = 4$ form with ${\tt sf\_trunc}$.
Bottom row: an exponential decline whose rise timescale $\tau$ is drawn independently
of $t_q$ (log-uniform on $[0.5, 5]$~Gyr) instead of being tied to it. Grey contours
enclose 68, 95 and 99.5 per cent of all epochs later than 1~Gyr; the red contour
encloses 68 and 95 per cent of the rapid-quenching class, which makes up 3.6, 6.4 and
0.84 per cent of the epochs in the three families (1.2 per cent for the exponential
family of Fig.~\ref{fig:index_planes_agb_on_off}), because a sharper cutoff holds the
recent-to-previous sSFR ratio below 0.1 for longer. In every family the rapid-quenching
class occupies the deepest part of the bump planes, with a median H$^-$ bump of
$-0.074$ to $-0.076$~mag, and the gain in completeness and purity from adding the bump
to D4000 and H$\delta_{\rm A}$ is positive with AGB on and consistent with zero with
AGB off.%
}
